\section{定性判据：无参数条件下的使用方式}

参数尚未标定时，框架先用于定位约束、进行有条件的比较和选择需要测量的环节。定性判断依靠公开证据与明确的序关系。能够判断的方向直接报告，存在权衡而证据不足的情形则保留为待比较的方案。

\subhead{定位：把对象放进三个坐标}

把一个对象放进 $(R,\,T,\,G)$ 三个坐标里，得到的不是一组数值，而是一个位置：它在三个维度上分别处在低、中、高哪一档。这一档的判断通常不需要测量，只需要回答表 \ref{tab:ch2-rtg} 列出的三个问题：推得动推不动、来得及来不及、传得开传不开。做完之后，对象就被分到了介质梯度的某一级上。

这种分类的操作性在于判据是公开的。两个人对同一个对象给出不同的档位，可以直接检查各自依据的是哪一条证据，而不必先争论一个数字谁更准。相比之下，“这项技术的黏滞系数是 3.7”这样的陈述无法被追究，因为没有人能说清 3.7 是从哪一步来的。模糊的档位在这件事上反而更硬：它逼着每一个判断落到一条可以指出来的证据上。

\subhead{分区：八个区域与各自的处理方式}

三个坐标各有 0 到 5 级，以中点 2.5 级为界，0 到 2 级记低，3 到 5 级记高。三个坐标各自或低或高，空间就被切成八个区域；判断一个对象落在哪一区，只需要回答三个是否：资源势垒是否过了中点，时间刚性是否过了中点，拓扑耦合是否过了中点。这个判断比定级更粗，因而也更稳，两个人即使在同一坐标上相差一级，只要没有跨过中点，得到的区域仍然相同。

区域的命名由高的坐标拼成。只有一个坐标高，就以它命名：资源高叫资源锁，时间高叫时间锁，耦合高叫耦合锁。两个坐标高，把两个名字连起来：资源与时间高叫资源时间，时间与耦合高叫时间耦合，资源与耦合高叫资源耦合。三个都高叫三重锁，三个都低叫低黏滞区。第六章按字面用过“低黏滞区”这个词，这里给它一个确切的边界。

\begin{figure}[ht]
\centering
% 原图自然宽度 13.40 cm，放大到 14.30 cm 以接近版心宽度
\scalebox{1.067}{% >>>>>>>>>> figs/fig-rtg-octant.tex >>>>>>>>>>
% 图 A-1  广义黏滞空间的三维分区与技术落位
%
% 投影：斜二测。R 向右、T 向上、G 向左上后退，均为线性映射
%   x = ox + L*(R - kx*G),  y = oy + L*(T + ky*G)
% 只画三个朝向读者的面：G=0 的前面、T=5 的顶面、R=0 的左面，共 12 个小
% 四边形。不可见的棱一律不画：唯一被挡住的顶点是 (5,0,5)（在图中投到
% (ox+1.368, 2.256)），从它引出的三条棱都在框体背后。把棱画成虚线会在
% 框内留下与中位线近于平行的第二条线（D 到 (5,0,5) 的棱与前面 T=2.5 的
% 中位线只差 0.14 cm），读者会误当成分区线，因此整体省略。
%
% 分区：三个坐标各以 2.5 级为界（即等级 0-2 记低、3-5 记高）。三个中位面
% 把空间切成八区。12 个小四边形按所属区域着色，深浅（色相）表示该区有
% 几个高坐标：绿 0、灰 1、棕 2、红 3。每个小四边形的名字由高坐标拼成，
% 例如 R 与 T 高称「资源时间」。
%
% 八个区域里有七个能在可见面上找到自己的格子：前面四格是 G 低的四区，
% 顶面四格是 T 高的四区，左面四格是 R 低的四区（四格重合于中位面两侧的
% 区）。剩下的一区是资源耦合（R 高、T 低、G 高），它朝向背面，没有可见
% 面，图中不画，由图注交代。
%
% 落位：每个技术点画在它自己的 (R,T,G) 上，再引一条细虚线连到它所属区域
% 在可见面上的落点（实心小点）。落点必定落在与其同区的格子里，于是「点
% 在哪一区」可以一步读出。规则：T 高就直上到 T=5，否则 G 低就沿 G 退到
% G=0，两者都不满足的点本身已经落在可见面上。
\begin{tikzpicture}[x=1cm, y=1cm, font=\small,
  ax/.style={-{Stealth[length=2.0mm]}, ink, line width=0.7pt},
  edge/.style={slate!75, line width=0.5pt},
  mid/.style={slate!55, line width=0.35pt, dash pattern=on 1.2pt off 1.1pt},
  stem/.style={line width=0.45pt, dash pattern=on 1.5pt off 1.1pt},
  foot/.style={circle, fill, inner sep=0.52pt},
  tick/.style={font=\tiny, text=slate, inner sep=0pt},
  rlab/.style={font=\scriptsize\heiti, inner sep=0pt},
  dlab/.style={font=\scriptsize\heiti, inner sep=0pt},
  clab/.style={font=\scriptsize\heiti, text=ink, inner sep=0pt},
  lead/.style={slate!75, line width=0.35pt},
  dom/.style={circle, inner sep=1.75pt},
  % 链上单元：实心圆，颜色取所属区域的深色
  ug/.style={circle, fill=forest, draw=white, line width=0.6pt, text=white,
             font=\tiny, minimum size=3.4mm, inner sep=0pt},
  ut/.style={circle, fill=terra, draw=white, line width=0.6pt, text=white,
             font=\tiny, minimum size=3.4mm, inner sep=0pt},
  ur/.style={circle, fill=rust, draw=white, line width=0.6pt, text=white,
             font=\tiny, minimum size=3.4mm, inner sep=0pt},
  um/.style={rounded corners=1.2pt, fill=forest, draw=white, line width=0.6pt,
             text=white, font=\tiny, inner sep=1.3pt},
  % 界面单元：空心圆，只有一圈彩色描边
  ii/.style={circle, fill=white, draw=terra, line width=0.8pt, text=terra,
             font=\tiny, minimum size=3.4mm, inner sep=0pt}
]
\useasboundingbox (-0.34,-0.72) rectangle (13.06,6.45);

% ================= 箱体：可见面着色、中位线、棱、坐标轴 =================
\newcommand{\boxat}[1]{%
  \coordinate (A) at ({#1},{0.60});
  \coordinate (B) at ({#1+3.60},{0.60});
  \coordinate (C) at ({#1},{4.20});
  \coordinate (D) at ({#1-2.232},{2.256});
  \coordinate (E) at ({#1+3.60},{4.20});
  \coordinate (G) at ({#1-2.232},{5.856});
  \coordinate (H) at ({#1+1.368},{5.856});
  % ---- 前面（G=0）切成四块 ----
  \fill[forest!16] (A) rectangle ({#1+1.80},{2.40});
  \fill[slate!17]  ({#1+1.80},{0.60}) rectangle ({#1+3.60},{2.40});
  \fill[slate!17]  ({#1},{2.40}) rectangle ({#1+1.80},{4.20});
  \fill[terra!19]  ({#1+1.80},{2.40}) rectangle (E);
  % ---- 左面（R=0）切成四块 ----
  \fill[forest!16] (A) -- ({#1},{2.40}) -- ({#1-1.116},{3.228}) -- ({#1-1.116},{1.428}) -- cycle;
  \fill[slate!17]  ({#1},{2.40}) -- (C) -- ({#1-1.116},{5.028}) -- ({#1-1.116},{3.228}) -- cycle;
  \fill[slate!17]  ({#1-1.116},{1.428}) -- ({#1-1.116},{3.228}) -- ({#1-2.232},{4.056}) -- ({#1-2.232},{2.256}) -- cycle;
  \fill[terra!19]  ({#1-1.116},{3.228}) -- ({#1-1.116},{5.028}) -- ({#1-2.232},{5.856}) -- ({#1-2.232},{4.056}) -- cycle;
  % ---- 顶面（T=5）切成四块 ----
  \fill[slate!17]  (C) -- ({#1+1.80},{4.20}) -- ({#1+0.684},{5.028}) -- ({#1-1.116},{5.028}) -- cycle;
  \fill[terra!19]  ({#1+1.80},{4.20}) -- (E) -- ({#1+2.484},{5.028}) -- ({#1+0.684},{5.028}) -- cycle;
  \fill[terra!19]  ({#1-1.116},{5.028}) -- ({#1+0.684},{5.028}) -- ({#1-0.432},{5.856}) -- ({#1-2.232},{5.856}) -- cycle;
  \fill[rust!17]   ({#1+0.684},{5.028}) -- ({#1+2.484},{5.028}) -- ({#1+1.368},{5.856}) -- ({#1-0.432},{5.856}) -- cycle;
  % ---- 可见面上的中位线 ----
  \draw[mid] ({#1+1.80},{0.60}) -- ({#1+1.80},{4.20});
  \draw[mid] ({#1},{2.40}) -- ({#1+3.60},{2.40});
  \draw[mid] ({#1-1.116},{1.428}) -- ({#1-1.116},{5.028});
  \draw[mid] ({#1},{2.40}) -- ({#1-2.232},{4.056});
  \draw[mid] ({#1+1.80},{4.20}) -- ({#1-0.432},{5.856});
  \draw[mid] ({#1-1.116},{5.028}) -- ({#1+2.484},{5.028});
  % ---- 可见的棱：外轮廓六边形与中间的三条 ----
  \draw[edge] (A) -- (B) -- (E) -- (C) -- cycle;
  \draw[edge] (E) -- (H) -- (G) -- (C);
  \draw[edge] (G) -- (D) -- (A);
  \draw[edge] (D) -- ({#1},{2.256});
  % ---- 坐标轴 ----
  \draw[ax] (A) -- ({#1+3.78},{0.60});
  \draw[ax] (A) -- ({#1},{4.38});
  \draw[ax] (A) -- ({#1-2.411},{2.389});
  \node[font=\scriptsize, text=ink, anchor=west] at ({#1+3.86},{0.60}) {$R$};
  \node[font=\scriptsize, text=ink, anchor=south east] at ({#1-0.05},{4.46}) {$T$};
  \node[font=\scriptsize, text=ink, anchor=east] at ({#1-2.50},{2.42}) {$G$};
}

% ================= (a) 八个区域与五个领域的落位 =================
\boxat{2.60}
\node[tick, anchor=north] at (2.50,0.48) {0};
\node[tick, anchor=north] at (6.20,0.48) {5};
\node[tick, anchor=east]  at (2.45,4.20) {5};
\node[tick, anchor=west]  at (0.44,2.08) {5};

% 区域名（颜色＝高坐标的个数；资源耦合没有可见面，由图注交代）
\node[rlab, text=forest] at (3.70,1.86) {低黏滞区};
\node[rlab, text=slate]  at (5.30,1.50) {资源锁};
\node[rlab, text=slate]  at (3.50,3.30) {时间锁};
\node[rlab, text=slate]  at (0.926,2.742) {耦合锁};
\node[rlab, text=terra]  at (5.30,3.30) {资源时间};
\node[rlab, text=terra]  at (0.926,4.542) {时间耦合};
\node[rlab, text=rust]   at (3.626,5.442) {三重锁};

% 五个领域：先画落位引线与落点，再画点，最后挂标签
\draw[stem, forest!80] (2.427,1.262) -- (3.32,0.60);
\node[foot, forest] at (3.32,0.60) {};
\draw[stem, rust!80] (2.974,4.805) -- (2.974,5.525);
\node[foot, rust] at (2.974,5.525) {};
\draw[stem, rust!80] (4.141,4.474) -- (4.141,5.194);
\node[foot, rust] at (4.141,5.194) {};
\node[dom, fill=forest] (d1) at (2.427,1.262) {};
\node[dom, fill=rust]   (d2) at (2.974,4.805) {};
\node[dom, fill=terra]  (d3) at (2.701,5.194) {};
\node[dom, fill=terra]  (d4) at (2.254,5.525) {};
\node[dom, fill=rust]   (d5) at (4.141,4.474) {};
\node[dlab, text=forest, anchor=west]  at (2.72,1.262) {大模型};
\node[dlab, text=rust,   anchor=west]  at (2.72,4.46)  {半导体};
\node[dlab, text=rust,   anchor=west]  at (4.32,4.40)  {脱碳};
\node[dlab, text=terra,  anchor=south] at (2.60,5.62)  {低空};
\node[dlab, text=terra,  anchor=south] at (1.66,5.90)  {基因编辑};
\draw[lead] (1.76,5.83) -- (2.19,5.56);
\draw[lead] (2.68,5.60) -- (2.70,5.30);

% ================= (b) 低空飞行器九个单元的落位 =================
\boxat{8.90}
\draw[stem, rust!80]   (9.721,3.754) -- (9.721,5.194);
\draw[stem, forest!80] (9.447,2.702) -- (10.34,2.04);
\draw[stem, forest!80] (9.174,1.651) -- (9.62,1.32);
\draw[stem, terra!80]  (9.001,3.754) -- (9.001,5.194);
\draw[stem, terra!80]  (7.834,4.805) -- (7.834,5.525);
\draw[stem, terra!80]  (7.388,5.136) -- (7.388,5.856);
\node[foot, rust]   at (9.721,5.194) {};
\node[foot, forest] at (10.34,2.04) {};
\node[foot, forest] at (9.62,1.32) {};
\node[foot, terra]  at (9.001,5.194) {};
\node[foot, terra]  at (7.834,5.525) {};
\node[foot, terra]  at (7.388,5.856) {};
\node[ur] at (9.721,3.754) {1};
\node[um] at (9.447,2.702) {2·3};
\node[ug] at (9.174,1.651) {4};
\node[ur] at (9.721,4.474) {5};
\node[ut] at (9.001,3.754) {6};
\node[ii] at (7.834,4.805) {7};
\node[ii] at (7.388,5.136) {8};
\node[ii] at (7.114,5.525) {9};
% 分类标签：实心＝链上，空心＝界面
\node[foot, ink, inner sep=0.62pt] at (9.28,3.16) {};
\node[clab, anchor=west] at (9.45,3.16) {链上单元};
\node[ii, minimum size=2.1mm, line width=0.9pt] at (7.02,6.14) {};
\node[clab, anchor=west] at (7.20,6.14) {界面单元};

% ================= 图例与栏标题（一律在内容下方）=================
\node[font=\scriptsize, text=slate, anchor=west] at (0.60,0.20) {高坐标的个数＝};
\foreach \i/\cc in {0/forest, 1/slate, 2/terra, 3/rust} {
  \pgfmathsetmacro{\gx}{3.10+0.90*\i}
  \fill[\cc] (\gx,0.11) rectangle ({\gx+0.20},0.31);
  \node[font=\scriptsize, text=slate, anchor=west] at ({\gx+0.30},0.21) {\i};
}
\node[anchor=north, font=\footnotesize\heiti, text=slate] at (3.28,-0.12)
      {(a) 八个区域与五个领域的落位};
\node[anchor=north, font=\footnotesize\heiti, text=slate] at (9.58,-0.12)
      {(b) 同一领域内部九个单元的落位};
\end{tikzpicture}

% <<<<<<<<<< figs/fig-rtg-octant.tex <<<<<<<<<<
}
\caption{三个坐标与八个区域。图中点位为分类方式的示意，展示单元与区域怎样对应；它们不构成实测行业评分、总体黏滞排序或投运时间预测。具体项目应按表\ref{tab:scale-five}的问题取证，并在同一判据下重新定位。}
\label{fig:rtg-octant}
\end{figure}

区域名中的“锁”表示当前任务面临的高约束，类别由所采用的量表确定。高约束是否可以放松，需要继续查明可替代资源、过程改造和结构重组的条件。分区先标出检查重点，处理方案再按证据选择。

低黏滞区的三个坐标均低，通常较少遇到这三类障碍。实际扩张仍取决于需求、可靠性和组织能力，分区本身不提供增长速度。

单锁区先检查主导约束：能否通过替代资源、改变过程或缩小依赖面来降低它。绕行是一种候选策略，其代价和可行性须与直接改造比较。

双锁区需要同时检查两项约束及其联动。两个坐标都高，不排除某项改进同时降低两者；改造节奏由可行路径、投入和验证要求决定。等值面斜率只描述给定模型中的代价组合。

三重锁区提示资源、时间和结构都需要重点核查。分阶段投入、重新定义任务或基础设施改造，都可能改变坐标组合。能否推进，需要检查项目的可行集合，不能由三项评分直接判定。

分区与定级是同一件事的两种精度。定级给出三个档次，分区给出一个格子；格子丢掉的是内部的差别，换来的是边界的稳定。代价也说清楚：靠近中点的对象会在两个区域之间摇摆，一次赋值的小改动就足以让它换区。这种摇摆不是缺陷，它对应现实中确实含混的情形，一项技术刚刚跨过中点时，说它“开始被卡住”和说它“还没被卡住”都不算错。

定位先定级、再定区，给出位置与需要优先核查的约束。下一节的量表把八个区域落到具体技术上。

\subhead{比较：共同改善与跨维权衡}

定位之后，两个对象的相对黏滞可以在不给权重数值的条件下比较出来。设两个状态 $a$、$b$ 在三个坐标上的差为 $\Delta R=R_{a}-R_{b}$、$\Delta T$、$\Delta G$，由内禀黏滞的指数形式两边取对数得

\begin{equation}
\ln\frac{\eta_{a}}{\eta_{b}}
=\alpha\,\Delta R+\beta\,\Delta T+\gamma\,\Delta G
\end{equation}

在参照与刻度一致、权重为正的条件下，若$a$的三个约束均不高于$b$，且至少一项更低，模型给出$\eta_a<\eta_b$。这是一种支配关系。若一项下降而另一项上升，只知道权重次序和变化符号通常不足以确定总效果，还需要权重比与变化幅度的界限。0到5级量表提供的是序数，代入指数模型之前，还须建立有意义的数量刻度。

\subhead{判方向：只需要导数的符号}

趋势判断需要的信息更少，只需要导数符号。对上式求时间导数，

\begin{equation}
\operatorname{sgn}\!\left(\frac{d\ln\eta}{dt}\right)
=\operatorname{sgn}\!\left(\alpha\dot{R}+\beta\dot{T}+\gamma\dot{G}\right)
\end{equation}

三个坐标同向变化且权重为正时，可以判断模型中的代价变化方向。不同方向同时出现，则需要比较加权幅度。系数随时间变化时，导数还包含$\dot\alpha R+\dot\beta T+\dot\gamma G$与基准变化项。生成成本下降与验证时间稳定，需要对各自环节测量，不能单由约束名称读出。

\subhead{模糊、开放、有弹性}

定性用法需要公开判据、记录不确定性并允许证据修正。

它是模糊的。档位之间的边界由判据给出，不由测量给出，因此一个对象可以同时符合相邻两档的描述。这种含混在证据不充分时，比一个精确但无根据的数字更接近实情。

判据与权重需要接受新证据。后续结果可能要求修改赋值、重新归一化、调整系数，也可能表明三维分类遗漏了机制。修订应按证据影响的范围进行，并留下对应记录。

判断可以随证据更新。一次赋值错误，需要复核相邻档位、跨区结论和依赖该赋值的推论；量化模型的参数也可以重新估计。修订范围由错误影响的链条决定。每一项判断都应能回到相应证据。

\subhead{限度}

定性用法回答不了三件事：多与少的差距有多大，快与慢的差距有多久，以及临界点落在哪里。上一节写下的 $r\theta$、正文用到的德博拉数与各章的阈值条件，都属于必须要有数值才能回答的问题，因此它们在本书里只以条件的形式出现，落不到具体对象上。

定性分析与定量分析的关系因此不是粗略与精确的关系，而是两组不同问题的关系：前者回答往哪个方向、在哪一档、谁更高，后者回答差多少、多久、何时越过。把前者当作后者的近似来用，会得到一个看起来有结论而实际无法检验的判断，这正是本书要避免的。反过来说，因为拿不到数值就放弃定性判断，同样是错的，那等于把三件本来能做对的事一起让掉。
